\appendix

\section{Proof Details}\label{sec:proof-details}

\subsection{Stable Selection Under Binary Recoding}\label{app:stable-selection}

\begin{lemma}[Exact measurable binary recoding]\label{lem:step-001-binary-recoding}
Let \(r\geq 1\), let \(C\subseteq\mathcal H_X\), let \(Q\) be a
distribution on \(Z_X=X\times\{0,1\}\), let \(h\in\mathcal H_X\), and
let \(G:Z_X^r\rightsquigarrow\mathcal H_X\) be a randomized kernel.
Then the pointwise bijection between the labels \(\{0,1\}\) and
\(\{-1,+1\}\), equipped with the transported output sigma-algebra,
preserves realizability, population zero-one risk, iid product laws,
exact output-atom probabilities, and ordered replacement adjacency in
both directions.
\end{lemma}

\begin{proof}
Define
\[
\lambda(0)=-1,\qquad \lambda(1)=+1,\qquad
\lambda^{-1}(t)=\frac{t+1}{2},
\]
and
\[
\Phi_Z(x,y)=(x,\lambda(y)),\qquad
\Phi_H(h)(x)=\lambda(h(x)).
\]
Both maps are pointwise bijections.  If \(\Sigma_X\) is the output
sigma-algebra on \(\mathcal H_X\), equip
\(\mathcal H_X^\pm=\{-1,+1\}^X\) with
\[
\Sigma_X^\pm=\{\Phi_H(E):E\in\Sigma_X\}.
\]
Because \(\Phi_H\) is bijective, \(\Sigma_X^\pm\) is a sigma-algebra
and both \(\Phi_H\) and \(\Phi_H^{-1}\) are measurable.  In particular,
measurable singleton hypotheses remain measurable.  On labeled
records we use the finite discrete label sigma-algebras and the
corresponding recordwise transport.  The standing measurability
convention therefore covers the equality, histogram, singleton, and
output events used below.

Set
\[
Q^\pm=(\Phi_Z)_\#Q,\qquad
G^\pm(S^\pm)=
\Phi_H\!\left(
G\bigl((\Phi_Z^{-1})^{\otimes r}(S^\pm)\bigr)
\right).
\tag{A.1}
\]
The transported structures make \(G^\pm\) a kernel.  If \(Q\) is
realizable by \(c\in C\), then for \(Q^\pm\)-almost every \((x,t)\),
\[
t=\lambda(c(x))=\Phi_H(c)(x),
\]
so \(Q^\pm\) is realizable by
\(\Phi_H(c)\in C^\pm:=\Phi_H(C)\).

For every \(h\in\mathcal H_X\), pointwise bijectivity gives
\[
\begin{aligned}
\operatorname{loss}_{Q^\pm}(\Phi_H(h))
&=\Pr_{(x,t)\sim Q^\pm}[\Phi_H(h)(x)\neq t]\\
&=\Pr_{(x,y)\sim Q}[\lambda(h(x))\neq\lambda(y)]\\
&=\Pr_{(x,y)\sim Q}[h(x)\neq y]
=R_Q(h).
\end{aligned}
\tag{A.2}
\]
The recordwise map gives the exact product-law identity
\[
(Q^\pm)^r=(\Phi_Z^{\otimes r})_\#Q^r.
\tag{A.3}
\]
Combining (A.1), (A.3), and bijectivity of \(\Phi_H\), including all
internal randomness of \(G\), yields
\[
\Pr_{S^\pm\sim(Q^\pm)^r,G^\pm}
[G^\pm(S^\pm)=\Phi_H(h)]
=
\Pr_{S\sim Q^r,G}[G(S)=h].
\tag{A.4}
\]

Finally, \(\Phi_Z\) is a bijection on individual records.  Hence for
ordered samples \(S,S'\) of the same length,
\[
\#\{i:S_i\neq S_i'\}
=
\#\{i:\Phi_Z(S_i)\neq\Phi_Z(S_i')\}.
\tag{A.5}
\]
Thus replacement of an instance, a label, or both components of one
record is carried to exactly one source-record replacement, and
conversely.  Equations (A.2)--(A.5) establish all asserted
zero-residual correspondences.
\end{proof}

\begin{theorem}[Stable selection from a frequent accurate hypothesis]
\label{thm:step-001-blm-selection}
Let \(\mathcal F\subseteq\{-1,+1\}^X\), let \(D\) be realizable by
\(\mathcal F\), let \(m\geq1\), let \(\eta\in(0,1]\), and let
\[
G:(X\times\{-1,+1\})^m
\rightsquigarrow\{-1,+1\}^X
\]
be randomized.  Suppose that some \(h\in\{-1,+1\}^X\) satisfies
\[
\Pr_{S\sim D^m,G}[G(S)=h]\geq\eta,
\qquad
\operatorname{loss}_D(h)\leq\frac{\alpha}{2}.
\]
For every \(\alpha,\beta,\varepsilon,\delta\in(0,1)\), there is an
arbitrary-output \((\varepsilon,\delta)\)-differentially private
kernel \(M:(X\times\{-1,+1\})^n
\rightsquigarrow\{-1,+1\}^X\) such that
\[
\Pr_{S\sim D^n,M}
[\operatorname{loss}_D(M(S))>\alpha]\leq\beta
\]
and
\[
n\leq C_{\mathrm{BLM}}
\left[
\frac{m}{\eta\varepsilon}
\log\!\left(\frac1{\eta\beta\delta}\right)
+
\frac1{\alpha\varepsilon}
\log\!\left(\frac1{\eta\beta}\right)
\right]
\tag{A.6}
\]
for a universal numerical constant \(C_{\mathrm{BLM}}>0\).
\end{theorem}

\begin{proof}[Source justification]
This is Theorem 17 in Section 5.2 of
\citet{bun2021equivalence}.  Its stable-histogram and private
selection construction partitions the input into disjoint
\(m\)-record batches, applies the fixed \(G\) to each batch, runs a
stable histogram on the resulting hypotheses, and then runs a generic
private learner on a fresh labeled subsample using the resulting
finite list.  The construction depends on
\(G,m,\eta,\alpha,\beta,\varepsilon,\delta\), but not on \(D\) or the
witness \(h\); the latter appears only in the utility analysis.  Its
codomain is the full \(\{-1,+1\}^X\), and its privacy convention is
symmetric fixed-length one-record replacement.  The source writes
(A.6) with absolute \(O(\cdot)\) constants.  A single
\(C_{\mathrm{BLM}}\) absorbs those universal implementation constants
and integer choices.  In the application below
\(\beta_0=1/8\), so both logarithms are at least \(\log 8\); any
additive integer unit is therefore absorbed without introducing a
parameter-dependent term.
\end{proof}

\begin{lemma}[Admissibility of stable selection]\label{lem:step-001-blm-admissibility}
Under Assumption~\ref{assump:polynomial-global-stability}, fix a
nonempty finite \(C\subseteq\{0,1\}^X\), put \(q=q(C)\) and
\(L=\log|C|\), and fix \(\varepsilon_0\in(0,1)\).  Let
\[
\eta_C=q^{-a},\qquad
\delta_C=\delta_a(q,L)=e^{-T_a(q,L)}.
\]
For every distribution \(Q\) realizable by \(C\), the recoded class
\(C^\pm\), the fixed producer \(G_C^\pm\), and the recoded witness
\(h_{C,Q}^\pm\) satisfy the hypotheses of
Theorem~\ref{thm:step-001-blm-selection} at
\[
(m,\eta,\alpha,\beta,\varepsilon,\delta)
=(m_C,\eta_C,\alpha_0,\beta_0,\varepsilon_0,\delta_C).
\]
If \(q=1\), then \(m_C=1\) and \(\eta_C=1\), and the same application
remains admissible.
\end{lemma}

\begin{proof}
Assumption~\ref{assump:polynomial-global-stability} supplies one
integer \(m_C\) and one randomized kernel \(G_C\), both fixed before
\(Q\), such that
\[
1\leq m_C\leq q^a.
\tag{A.7}
\]
For every realizable \(Q\), it supplies an arbitrary
\(h_{C,Q}\in\mathcal H_X\) satisfying
\[
R_Q(h_{C,Q})\leq\frac{\alpha_0}{2},
\qquad
\Pr_{S\sim Q^{m_C},G_C}[G_C(S)=h_{C,Q}]
\geq q^{-a}=\eta_C.
\tag{A.8}
\]
Apply Lemma~\ref{lem:step-001-binary-recoding} with \(r=m_C\).
Equations (A.2) and (A.4) turn (A.8) into
\[
\operatorname{loss}_{Q^\pm}(h_{C,Q}^\pm)
\leq\frac{\alpha_0}{2},
\qquad
\Pr[G_C^\pm((Q^\pm)^{m_C})=h_{C,Q}^\pm]
\geq\eta_C.
\tag{A.9}
\]
The same lemma proves that \(Q^\pm\) is realizable by \(C^\pm\).
Neither the witness nor the producer output is required to lie in
\(C^\pm\).

Since \(q\geq1\) and \(a\geq1\), \(0<\eta_C\leq1\).
The values \(\alpha_0,\beta_0,\varepsilon_0\) lie in \((0,1)\).
Moreover, \(q^b\geq1\) and \(u(L)\geq1\), so
\[
T_a(q,L)\geq2,\qquad
0<\delta_C=e^{-T_a(q,L)}\leq e^{-2}<1.
\tag{A.10}
\]
At \(q=1\), (A.7) forces \(m_C=1\), while
\(\eta_C=1^{-a}=1\).  The logarithms and denominators in (A.6) remain
legal because \(\beta_0,\delta_C\in(0,1)\).  Thus no \(q>1\)
restriction has been introduced.
\end{proof}

\begin{proposition}[A distribution-free source-label learner]
\label{prop:step-001-uniform-blm}
Under Assumption~\ref{assump:polynomial-global-stability},
Lemma~\ref{lem:step-001-binary-recoding},
Lemma~\ref{lem:step-001-blm-admissibility}, and
Theorem~\ref{thm:step-001-blm-selection}, fix a nonempty finite \(C\),
put \(q=q(C)\), \(L=\log|C|\), and fix
\(\varepsilon_0\in(0,1)\).  Then there are one integer \(n_C\) and one
kernel
\[
A_C^\pm:(X\times\{-1,+1\})^{n_C}
\rightsquigarrow\{-1,+1\}^X,
\]
both independent of the realizable distribution, such that for every
realizable \(Q\),
\[
\Pr_{S^\pm\sim(Q^\pm)^{n_C},A_C^\pm}
\!\left[
\operatorname{loss}_{Q^\pm}(A_C^\pm(S^\pm))>\alpha_0
\right]\leq\beta_0,
\tag{A.11}
\]
\(A_C^\pm\) is \((\varepsilon_0,\delta_C)\)-differentially private,
and
\[
n_C\leq C_{\mathrm{BLM}}
\left[
\frac{m_C}{q^{-a}\varepsilon_0}
\log\!\left(\frac1{q^{-a}\beta_0\delta_C}\right)
+
\frac1{\alpha_0\varepsilon_0}
\log\!\left(\frac1{q^{-a}\beta_0}\right)
\right].
\tag{A.12}
\]
The output range is all of \(\{-1,+1\}^X\).
\end{proposition}

\begin{proof}
Fix \(C,G_C^\pm,m_C\), and the public parameter tuple from
Lemma~\ref{lem:step-001-blm-admissibility}.  Fix also the
stable-histogram and private-selection implementations from
Theorem~\ref{thm:step-001-blm-selection}.  Their construction invokes
only the fixed \(G_C^\pm\) and the public tuple
\[
(m_C,q^{-a},\alpha_0,\beta_0,\varepsilon_0,\delta_C).
\]
Neither \(Q^\pm\) nor \(h_{C,Q}^\pm\) is an algorithm input or
implementation choice.  Thus the prescription defines one kernel
\(A_C^\pm\) and one integer \(n_C\), independent of \(Q\).
For each realizable \(Q\),
Lemma~\ref{lem:step-001-blm-admissibility} supplies a witness for that
same fixed kernel.  Applying its utility conclusion
separately to each \(Q\) gives (A.11), with the quantifier order
\[
\exists A_C^\pm\quad
\forall Q\text{ realizable by }C.
\]

Its privacy proof is distribution-independent and supplies
both ordered replacement inequalities for every adjacent input pair.
Instantiating (A.6) gives (A.12), with the same universal
\(C_{\mathrm{BLM}}\) for all \(X,C,Q,m_C,q,L\) and all admissible
\(\varepsilon_0\).  Both terms remain intact: this argument has not
replaced \(m_C\) by \(q^a\), expanded the logarithms, dominated the
\(\alpha_0^{-1}\) term, selected \(K_a\), or invoked a sequence
limit.
\end{proof}

\begin{proposition}[Decoded arbitrary-output private learner]
\label{prop:step-001-binary-private-learner}
Under Assumption~\ref{assump:polynomial-global-stability},
Lemma~\ref{lem:step-001-binary-recoding}, and
Proposition~\ref{prop:step-001-uniform-blm}, define, for
\(S\in Z_X^{n_C}\),
\[
A_C(S)=
\Phi_H^{-1}\!\left(
A_C^\pm(\Phi_Z^{\otimes n_C}(S))
\right).
\tag{A.13}
\]
Then \(A_C:Z_X^{n_C}\rightsquigarrow\mathcal H_X\) is one
arbitrary-output, distribution-free realizable
\((\alpha_0,\beta_0)\)-PAC learner for \(C\), is
\((\varepsilon_0,\delta_C)\)-differentially private under ordered
replacement adjacency, and satisfies the sample bound (A.12).
\end{proposition}

\begin{proof}
The maps in (A.13) are bimeasurable by
Lemma~\ref{lem:step-001-binary-recoding}, so \(A_C\) is a kernel fixed
independently of \(Q\).  If \(Q\) is realizable by \(C\) and
\(S\sim Q^{n_C}\), then
\(\Phi_Z^{\otimes n_C}(S)\sim(Q^\pm)^{n_C}\).  For every source output
\(g^\pm\), (A.2) gives
\[
R_Q(\Phi_H^{-1}(g^\pm))
=\operatorname{loss}_{Q^\pm}(g^\pm).
\]
Consequently (A.11) implies
\[
\Pr_{S\sim Q^{n_C},A_C}
[R_Q(A_C(S))>\alpha_0]\leq\beta_0.
\tag{A.14}
\]
The same \(A_C\) satisfies (A.14) for every realizable \(Q\), and the
probability includes both sample and learner randomness.

For privacy, let \(S,S'\in Z_X^{n_C}\) be replacement-adjacent and let
\(E\subseteq\mathcal H_X\) be measurable.  By
Lemma~\ref{lem:step-001-binary-recoding},
\(\Phi_Z^{\otimes n_C}(S)\) and
\(\Phi_Z^{\otimes n_C}(S')\) are source-adjacent and
\(E^\pm:=\Phi_H(E)\) is measurable.  Hence
\[
\begin{aligned}
\Pr[A_C(S)\in E]
&=\Pr[A_C^\pm(\Phi_Z^{\otimes n_C}(S))\in E^\pm]\\
&\leq e^{\varepsilon_0}
\Pr[A_C^\pm(\Phi_Z^{\otimes n_C}(S'))\in E^\pm]
+\delta_C\\
&=e^{\varepsilon_0}\Pr[A_C(S')\in E]+\delta_C.
\end{aligned}
\tag{A.15}
\]
Interchanging \(S,S'\) gives the reverse inequality.  Identical
samples are immediate, while (A.5) covers replacement of the
instance, label, or both components of one record.

The inverse output map is deterministic postprocessing whose image is
all of \(\mathcal H_X\).  Thus arbitrary improper outputs and the lack
of representation or computational restrictions are preserved, and
the sample bound (A.12) is unchanged.  Together,
Lemma~\ref{lem:step-001-binary-recoding},
Lemma~\ref{lem:step-001-blm-admissibility},
Proposition~\ref{prop:step-001-uniform-blm}, and the present argument
derive the complete stable-selection conclusion without invoking any
later asymptotic specialization.
\end{proof}
\subsection{Explicit Sample Domination and Exact Padding}
\label{app:sample-domination}

\begin{lemma}[Explicit stable-selection specialization]
\label{lem:step-002-sc-domination}
Under Assumption~\ref{assump:polynomial-global-stability} and
Proposition~\ref{prop:step-001-binary-private-learner}, fix a
nonempty finite \(C\), put \(q=q(C)\), \(L=\log|C|\), and fix
\(\varepsilon_0\in(0,1)\).  Define
\[
B_0=\log\frac1{\beta_0}=\log 8,
\qquad
D_a=1+a+B_0+\alpha_0^{-1}(a+B_0)
=1+9a+9\log 8,
\tag{A.16}
\]
and choose
\[
K_a=\max\{2,C_{\mathrm{BLM}}D_a\}.
\tag{A.17}
\]
If \(n_C\) is the sample size in
Proposition~\ref{prop:step-001-binary-private-learner}, then
\[
\begin{aligned}
n_C
&\leq \frac{C_{\mathrm{BLM}}}{\varepsilon_0}
\left[
q^{2a}\!\left(
a\log q+\log\frac1{\beta_0}+T_a(q,L)
\right)
+\frac1{\alpha_0}
\left(a\log q+\log\frac1{\beta_0}\right)
\right]\\
&\leq K_a\varepsilon_0^{-2}q^b(1+T_a(q,L))
\leq N_a(q,L,\varepsilon_0).
\end{aligned}
\tag{A.18}
\]
The constant \(K_a\) depends only on
\(a,\alpha_0,\beta_0,C_{\mathrm{BLM}}\), and the inequalities remain
valid at \(q=1\) and uniformly over \(0<\varepsilon_0<1\).
\end{lemma}

\begin{proof}
Write \(T=T_a(q,L)\) and
\(\delta_C=\delta_a(q,L)=e^{-T}\).  From
Assumption~\ref{assump:polynomial-global-stability},
\[
\frac{m_C}{q^{-a}}=m_Cq^a\leq q^{2a}.
\tag{A.19}
\]
The two logarithms in (A.12) expand exactly as
\[
\log\!\left(\frac1{q^{-a}\beta_0\delta_C}\right)
=\log\!\left(\frac{q^ae^T}{\beta_0}\right)
=a\log q+B_0+T,
\tag{A.20}
\]
\[
\log\!\left(\frac1{q^{-a}\beta_0}\right)
=\log\!\left(\frac{q^a}{\beta_0}\right)
=a\log q+B_0.
\tag{A.21}
\]
Substitution of (A.19)--(A.21) into (A.12) gives the first inequality
of (A.18), with every confidence and logarithmic term retained.

Since \(q\geq1\), \(b=2a+2\), and
\(T=q^b+u(L)\),
\[
q^{2a}\leq q^b\leq T,\qquad
a\log q\leq aq^b\leq aT,\qquad
q^b\geq1.
\tag{A.22}
\]
The five summands are controlled separately:
\[
q^{2a}T\leq q^bT\leq q^b(1+T),
\tag{A.23}
\]
\[
aq^{2a}\log q
\leq aq^{2a}T
\leq aq^bT
\leq aq^b(1+T),
\tag{A.24}
\]
\[
B_0q^{2a}
\leq B_0q^b
\leq B_0q^b(1+T),
\tag{A.25}
\]
\[
\alpha_0^{-1}a\log q
\leq\alpha_0^{-1}aq^b
\leq\alpha_0^{-1}aq^b(1+T),
\tag{A.26}
\]
and
\[
\alpha_0^{-1}B_0
\leq\alpha_0^{-1}B_0q^b
\leq\alpha_0^{-1}B_0q^b(1+T).
\tag{A.27}
\]
Thus
\[
q^{2a}(T+a\log q+B_0)
+\alpha_0^{-1}(a\log q+B_0)
\leq D_aq^b(1+T).
\tag{A.28}
\]
Using (A.17) and
\(\varepsilon_0^{-1}\leq\varepsilon_0^{-2}\) gives
\[
\frac{C_{\mathrm{BLM}}}{\varepsilon_0}
D_aq^b(1+T)
\leq
K_a\varepsilon_0^{-2}q^b(1+T).
\tag{A.29}
\]
The integer \(n_C\) is at most this real quantity and hence at most
its ceiling \(N_a(q,L,\varepsilon_0)\).

At \(q=1\), one has \(m_C=1\), \(q^{-a}=1\),
\(\log q=0\), \(q^{2a}=q^b=1\), and \(T\geq1\), so
(A.19)--(A.29) remain literal.  No division by \(\log q\) or
assumption \(q>1\) occurs.  The inequality
\(\varepsilon_0^{-1}\leq\varepsilon_0^{-2}\) holds throughout
\((0,1)\) and tends to equality as \(\varepsilon_0\uparrow1\), so
\(K_a\) has no dependence on a gap from \(1\).
\end{proof}

\begin{proposition}[Exact-size padding preserves utility and privacy]
\label{prop:step-002-exact-padding}
Under Assumption~\ref{assump:polynomial-global-stability},
Proposition~\ref{prop:step-001-binary-private-learner}, and
Lemma~\ref{lem:step-002-sc-domination}, put
\[
N=N_a(q(C),\log|C|,\varepsilon_0).
\]
Then there is a kernel
\[
M_C:Z_X^N\rightsquigarrow\mathcal H_X
\]
with exact input arity \(N\), arbitrary output, distribution-free
realizable \((\alpha_0,\beta_0)\)-PAC utility, and
\((\varepsilon_0,\delta_a(q(C),\log|C|))\)-differential privacy.
No properness, representation, or computational restriction is
introduced.
\end{proposition}

\begin{proof}
Lemma~\ref{lem:step-002-sc-domination} gives \(n_C\leq N\).  Let
\[
\pi_{n_C}^N(z_1,\ldots,z_N)=(z_1,\ldots,z_{n_C})
\tag{A.30}
\]
and define
\[
M_C=A_C\circ\pi_{n_C}^N.
\tag{A.31}
\]
Thus \(M_C\) accepts exactly \(N\) records and ignores the final
\(N-n_C\).  If \(N=n_C\), it equals \(A_C\).

For every distribution \(Q\) realizable by \(C\), the prefix of a
sample \(S\sim Q^N\) has law \(Q^{n_C}\).  Therefore the full output
laws, including internal randomness, agree and
\[
\begin{aligned}
\Pr_{S\sim Q^N,M_C}[R_Q(M_C(S))>\alpha_0]
&=
\Pr_{S_0\sim Q^{n_C},A_C}
[R_Q(A_C(S_0))>\alpha_0]\\
&\leq\beta_0.
\end{aligned}
\tag{A.32}
\]
The kernel, its arity, and the public projection are fixed before
\(Q\), so this is distribution-free utility.

If \(S,S'\in Z_X^N\) are replacement-adjacent, projection cannot
increase the number of differing records:
\[
\#\{i\leq n_C:S_i\neq S_i'\}
\leq
\#\{i\leq N:S_i\neq S_i'\}
\leq1.
\tag{A.33}
\]
For every measurable \(E\subseteq\mathcal H_X\),
\[
\begin{aligned}
\Pr[M_C(S)\in E]
&=\Pr[A_C(\pi_{n_C}^N(S))\in E]\\
&\leq e^{\varepsilon_0}
\Pr[A_C(\pi_{n_C}^N(S'))\in E]
+\delta_a(q(C),\log|C|)\\
&=e^{\varepsilon_0}\Pr[M_C(S')\in E]
+\delta_a(q(C),\log|C|).
\end{aligned}
\tag{A.34}
\]
Interchanging \(S,S'\) gives the reverse inequality.  If the changed
record is in the ignored suffix, the two prefixes and hence the two
output laws are identical; if it lies in the prefix, (A.33) invokes
the original replacement comparison.  Padding therefore has zero
utility and privacy residual.  Combining (A.18) with this projection
establishes the exact sample-size and learner interface used below.
\end{proof}

\subsection{The Exact Privacy Schedule}\label{app:privacy-schedule}

\begin{lemma}[The cardinality correction]\label{lem:step-003-log-correction}
For every \(L\geq0\), the quantity
\[
u(L)=\log\log(e^e+L)
\]
is well-defined and satisfies
\[
u(L)\geq1,\qquad
e^{-u(L)}=\frac1{\log(e^e+L)}.
\tag{A.35}
\]
For every \(R>0\),
\[
L\geq\exp(\exp R)\quad\Longrightarrow\quad u(L)\geq R.
\tag{A.36}
\]
Consequently \(u(L)\to\infty\) as \(L\to\infty\).  At \(L=0\),
\(u(0)=1\) and \(e^{-u(0)}=e^{-1}\).
\end{lemma}

\begin{proof}
For \(L\geq0\),
\[
e^e+L\geq e^e>1,\qquad
\log(e^e+L)\geq\log(e^e)=e>0.
\tag{A.37}
\]
Thus both logarithms are defined, and monotonicity gives
\[
u(L)=\log\log(e^e+L)\geq\log e=1.
\tag{A.38}
\]
Exact exponentiation yields
\[
e^{-u(L)}
=\exp[-\log\log(e^e+L)]
=\frac1{\log(e^e+L)}.
\tag{A.39}
\]
If \(L\geq\exp(\exp R)\), then
\[
u(L)\geq
\log\log(\exp(\exp R))=R,
\tag{A.40}
\]
which is the quantified divergence statement.  Finally,
\[
u(0)=\log\log(e^e)=1,\qquad
e^{-u(0)}=e^{-1}=\frac1{\log(e^e)}.
\tag{A.41}
\]
\end{proof}

\begin{proposition}[Exact privacy logarithm]\label{prop:step-003-exact-privacy-schedule}
Under Assumption~\ref{assump:polynomial-global-stability} and
Lemma~\ref{lem:step-003-log-correction}, fix \(a\geq1\), put
\(b=2a+2\), and let \(q\geq1\), \(L\geq0\).  Then
\[
\delta_a(q,L)
=\exp[-q^b-\log\log(e^e+L)]
=\frac{e^{-q^b}}{\log(e^e+L)}
\in(0,e^{-2}],
\tag{A.42}
\]
and
\[
\log\frac1{\delta_a(q,L)}
=T_a(q,L)
=q^{2a+2}+\log\log(e^e+L).
\tag{A.43}
\]
For every finite \(C\), the learner in
Proposition~\ref{prop:step-002-exact-padding} is private at precisely
this parameter.  At \(L=0\) the privacy logarithm is \(q^b+1\); at
\(q=1\) it is \(1+\log\log(e^e+L)\); and at \((q,L)=(1,0)\) it is
\(2\).
\end{proposition}

\begin{proof}
Since \(a\geq1\),
\[
b=2a+2\geq4>0.
\tag{A.44}
\]
Thus \(q^b\geq1\), while
Lemma~\ref{lem:step-003-log-correction} gives \(u(L)\geq1\).  Hence
\[
T_a(q,L)=q^b+u(L)\geq2
\tag{A.45}
\]
and
\[
0<\delta_a(q,L)=e^{-T_a(q,L)}\leq e^{-2}<1.
\tag{A.46}
\]
Using (A.39),
\[
\delta_a(q,L)
=e^{-q^b}e^{-u(L)}
=\frac{e^{-q^b}}{\log(e^e+L)},
\tag{A.47}
\]
which proves (A.42).  Positivity permits reciprocation, so
\[
\begin{aligned}
\log\frac1{\delta_a(q,L)}
&=\log(e^{T_a(q,L)})\\
&=T_a(q,L)
=q^b+u(L)\\
&=q^{2a+2}+\log\log(e^e+L).
\end{aligned}
\tag{A.48}
\]
This is an equality, not an asymptotic simplification.

Proposition~\ref{prop:step-002-exact-padding} gives
\((\varepsilon_0,\delta_a(q,L))\)-privacy.  Substituting the identical
quantity from (A.42) changes neither the kernel nor the privacy
allowance.  At \(L=0\),
\[
\delta_a(q,0)=e^{-q^b-1},\qquad
\log\frac1{\delta_a(q,0)}=q^b+1.
\tag{A.49}
\]
At \(q=1\),
\[
\delta_a(1,L)=\frac{e^{-1}}{\log(e^e+L)},\qquad
\log\frac1{\delta_a(1,L)}
=1+\log\log(e^e+L).
\tag{A.50}
\]
These yield \(\delta_a(1,0)=e^{-2}\), with no requirement that
\(q>1\) or \(L>0\).
\end{proof}

\begin{proposition}[Activation on bounded-complexity sequences]
\label{prop:step-003-bounded-q-activation}
Under Proposition~\ref{prop:step-003-exact-privacy-schedule}, let
\((q_\kappa,L_\kappa)\) satisfy
\[
\begin{aligned}
&L_\kappa\geq0,\qquad L_\kappa\to\infty,\\
&1\leq q_\kappa\leq\bar q
\end{aligned}
\tag{A.51}
\]
for a finite \(\bar q\geq1\).  Then
\[
\delta_a(q_\kappa,L_\kappa)
=\frac{e^{-q_\kappa^b}}{\log(e^e+L_\kappa)}
\longrightarrow0,
\tag{A.52}
\]
even though \(q_\kappa\) need not converge.  More precisely,
\[
e^{-\bar q^b}\leq e^{-q_\kappa^b}\leq e^{-1},
\tag{A.53}
\]
so the cardinality factor, rather than the polynomial-in-\(q_\kappa\)
factor, forces the decay.  In particular,
\[
\delta_a(1,L_\kappa)
=\frac{e^{-1}}{\log(e^e+L_\kappa)}
\longrightarrow0.
\tag{A.54}
\]
\end{proposition}

\begin{proof}
From (A.51) and \(b>0\),
\(1\leq q_\kappa^b\leq\bar q^b\).  Since \(x\mapsto e^{-x}\) is
decreasing, this proves (A.53).  Thus a schedule containing only
\(e^{-q_\kappa^b}\) would stay at least
\(e^{-\bar q^b}>0\) on this sequence.

In contrast, Lemma~\ref{lem:step-003-log-correction} gives
\[
u(L_\kappa)\to\infty,\qquad
e^{-u(L_\kappa)}
=\frac1{\log(e^e+L_\kappa)}
\longrightarrow0.
\tag{A.55}
\]
Using (A.47) and the upper bound in (A.53),
\[
0<\delta_a(q_\kappa,L_\kappa)
\leq\frac{e^{-1}}{\log(e^e+L_\kappa)}
\longrightarrow0.
\tag{A.56}
\]
Equation (A.54) is the exact specialization \(q_\kappa=1\).
Accordingly, the explicit cardinality correction preserves privacy
decay in the bounded-\(q_\kappa\) regime without assuming
\(q_\kappa\to\infty\).
\end{proof}

\subsection{Subpower Consequences of Superpolynomial Separation}
\label{app:subpower}

Throughout this subsection, let \(\{C_\kappa\}\) be a sequence of
nonempty finite classes satisfying \(|C_\kappa|\to\infty\), the
counting bound
\[
0\leq\ell_\kappa\leq\log_2|C_\kappa|
=\frac{L_\kappa}{\log2},
\]
and
\[
\forall p\in\mathbb N\ \exists\kappa_0(p)\
\forall\kappa\geq\kappa_0(p):
\qquad L_\kappa>d_\kappa^p.
\tag{A.57}
\]
Then \(L_\kappa=\log|C_\kappa|\to\infty\), so all subsequent
eventual bounds may be restricted to \(L_\kappa\geq1\).

\begin{lemma}[Every-exponent control of the iterated-log term]
\label{lem:step-004-iterated-log-subpower}
Under the conditions above, for every \(r>0\) there is
\(\kappa_s(r)\) such that
\[
\kappa\geq\kappa_s(r)
\quad\Longrightarrow\quad
s_\kappa=1+\log^*(1+\ell_\kappa)
\leq L_\kappa^r.
\tag{A.58}
\]
This includes \(\ell_\kappa=0\), for which \(s_\kappa=1\).
\end{lemma}

\begin{proof}
The stopping time \(x\mapsto\log^*x\) is nondecreasing on
\([1,\infty)\).  Indeed, if \(1\leq x\leq y\) and
\(m=\log^*y\), then either the iterates of \(x\) reach a value at most
\(1\) before time \(m\), or monotonicity of \(\log_2\) through the
active iterates gives
\(\log_2^{(m)}x\leq\log_2^{(m)}y\leq1\).

For each integer \(n\geq1\),
\[
\log^*(2^n)=1+\log^*n\leq n.
\tag{A.59}
\]
The claim holds for \(n=1\).  If it holds at \(n-1\) and \(n\geq2\),
then \(n\leq2^{n-1}\), so monotonicity gives
\(\log^*n\leq\log^*(2^{n-1})\leq n-1\).  For \(x>1\), taking
\(n=\lceil\log_2x\rceil\) yields
\[
\log^*x
\leq n
\leq1+\log_2x.
\tag{A.60}
\]
At \(x=1\), \(\log^*1=0\), so the same inequality holds.

When \(L_\kappa\geq1\),
\[
1+\ell_\kappa
\leq1+\frac{L_\kappa}{\log2}
\leq c_2L_\kappa,
\qquad
c_2=1+\frac1{\log2}.
\tag{A.61}
\]
Using (A.60),
\[
\begin{aligned}
s_\kappa
&\leq2+\log_2(1+\ell_\kappa)\\
&\leq2+\log_2(c_2L_\kappa)
=A_0+\frac{\log L_\kappa}{\log2},
\end{aligned}
\qquad
A_0=2+\frac{\log c_2}{\log2}.
\tag{A.62}
\]
Fix \(r>0\).  For \(L\geq1\),
\[
\log L
=\int_1^L\frac{dx}{x}
\leq\int_1^Lx^{r/2-1}\,dx
=\frac2r(L^{r/2}-1)
\leq\frac2rL^{r/2}.
\tag{A.63}
\]
Thus
\[
s_\kappa\leq
A_rL_\kappa^{r/2},
\qquad
A_r=A_0+\frac{2}{r\log2}.
\tag{A.64}
\]
Since \(L_\kappa\to\infty\), eventually
\(L_\kappa^{r/2}\geq A_r\), and then
\(A_rL_\kappa^{r/2}\leq L_\kappa^r\).  If
\(\ell_\kappa=0\), then \(\log^*1=0\) and the same eventual bound
holds without any division by \(\ell_\kappa\).
\end{proof}

\begin{proposition}[Superpolynomial separation implies subpower total complexity]
\label{prop:step-004-q-subpower}
Under (A.57), the size-divergence and counting conditions, and
Lemma~\ref{lem:step-004-iterated-log-subpower},
\[
\forall t>0\ \exists\kappa_q(t)\
\forall\kappa\geq\kappa_q(t):
\qquad q_\kappa\leq L_\kappa^t.
\tag{A.65}
\]
Consequently \(q_\kappa=L_\kappa^{o(1)}\).  The conclusion permits
\(d_\kappa=0\) and bounded or nonconvergent \(q_\kappa\).
\end{proposition}

\begin{proof}
Fix \(t>0\) and choose one integer
\[
p=p(t)>\frac2t,\qquad \frac1p<\frac t2.
\tag{A.66}
\]
Condition (A.57) gives, for every sufficiently large \(\kappa\),
\[
L_\kappa>d_\kappa^p.
\tag{A.67}
\]
After also requiring \(L_\kappa\geq1\), taking the positive
\(p\)-th root gives
\[
d_\kappa<L_\kappa^{1/p}
\leq L_\kappa^{t/2}.
\tag{A.68}
\]
This remains valid at \(d_\kappa=0\).  Applying
Lemma~\ref{lem:step-004-iterated-log-subpower} with exponent \(t/2\)
gives
\[
s_\kappa\leq L_\kappa^{t/2}
\tag{A.69}
\]
on another tail.  After additionally requiring
\(L_\kappa^{t/2}\geq2\),
\[
q_\kappa=d_\kappa+s_\kappa
<2L_\kappa^{t/2}
\leq L_\kappa^t.
\tag{A.70}
\]
The threshold is the maximum of finitely many thresholds and
therefore applies to every later index.  Since \(q_\kappa\geq1\) and
eventually \(L_\kappa>1\), (A.65) also gives
\[
0\leq
\frac{\log q_\kappa}{\log L_\kappa}
\leq t
\quad\text{eventually}.
\tag{A.71}
\]
Letting the arbitrary \(t\) decrease to zero proves the
\(L_\kappa^{o(1)}\) statement.
\end{proof}

\begin{lemma}[Subpower growth of the privacy schedule]
\label{lem:step-004-schedule-subpower}
Under Assumption~\ref{assump:polynomial-global-stability},
Lemma~\ref{lem:step-003-log-correction},
Proposition~\ref{prop:step-003-exact-privacy-schedule}, and
Proposition~\ref{prop:step-004-q-subpower}, for every \(r>0\) there is
\(\kappa_T(r)\) such that
\[
\kappa\geq\kappa_T(r)
\quad\Longrightarrow\quad
u(L_\kappa)\leq L_\kappa^r,\qquad
T_a(q_\kappa,L_\kappa)\leq L_\kappa^r.
\tag{A.72}
\]
Hence \(u(L_\kappa)=L_\kappa^{o(1)}\) and
\(T_a(q_\kappa,L_\kappa)=L_\kappa^{o(1)}\), including when
\(q_\kappa\) is bounded.
\end{lemma}

\begin{proof}
For \(L\geq1\), put \(z=\log(e^e+L)>1\).  Then
\[
\log z=\int_1^z\frac{dx}{x}\leq z-1\leq z.
\]
Also \(e^e+L\leq(e^e+1)L\), so
\[
u(L)
\leq\log(e^e+L)
\leq\log(e^e+1)+\log L.
\tag{A.73}
\]
Fix \(r>0\).  Applying (A.63) gives
\[
u(L_\kappa)
\leq B_rL_\kappa^{r/2},
\qquad
B_r=\log(e^e+1)+\frac2r.
\tag{A.74}
\]
Eventually \(L_\kappa^{r/2}\geq B_r\), hence
\[
u(L_\kappa)\leq L_\kappa^r.
\tag{A.75}
\]

Now fix a fresh target exponent \(r>0\).
Proposition~\ref{prop:step-004-q-subpower}, used at
\(r/(2b)>0\), gives
\[
q_\kappa^b\leq L_\kappa^{r/2}
\tag{A.76}
\]
eventually.  The first part, used at exponent \(r/2\), gives
\[
u(L_\kappa)\leq L_\kappa^{r/2}
\tag{A.77}
\]
eventually.  On a common tail with \(L_\kappa^{r/2}\geq2\),
\[
T_a(q_\kappa,L_\kappa)
=q_\kappa^b+u(L_\kappa)
\leq2L_\kappa^{r/2}
\leq L_\kappa^r.
\tag{A.78}
\]
The lower bounds \(u(L_\kappa)\geq1\) and
\(T_a(q_\kappa,L_\kappa)\geq2\), together with (A.72), yield
\[
\frac{\log u(L_\kappa)}{\log L_\kappa}\to0,\qquad
\frac{\log T_a(q_\kappa,L_\kappa)}{\log L_\kappa}\to0.
\tag{A.79}
\]
If \(q_\kappa\leq\bar q\) on a tail, then directly
\[
T_a(q_\kappa,L_\kappa)
\leq\bar q^b+u(L_\kappa),
\tag{A.80}
\]
and the fixed first term and (A.75) are absorbed by every positive
power of \(L_\kappa\).
\end{proof}

\begin{proposition}[The exact sample size is subpower and sublogarithmic]
\label{prop:step-004-sample-subpower}
Under Assumption~\ref{assump:polynomial-global-stability},
Lemma~\ref{lem:step-002-sc-domination},
Proposition~\ref{prop:step-004-q-subpower}, and
Lemma~\ref{lem:step-004-schedule-subpower}, let
\[
N_\kappa=N_a(q_\kappa,L_\kappa,\varepsilon_0).
\]
Then
\[
\forall r>0\ \exists\kappa_N(r)\
\forall\kappa\geq\kappa_N(r):
\qquad N_\kappa\leq L_\kappa^r.
\tag{A.81}
\]
Consequently
\[
N_\kappa=L_\kappa^{o(1)},\qquad
N_\kappa=o(L_\kappa).
\tag{A.82}
\]
Both conclusions remain true for bounded \(q_\kappa\), including
\(q_\kappa\equiv1\).
\end{proposition}

\begin{proof}
Set
\[
H=K_a\varepsilon_0^{-2}.
\tag{A.83}
\]
By Lemma~\ref{lem:step-002-sc-domination}, \(H\) is finite and fixed
while \(\kappa\to\infty\), with dependence only on
\(a,\alpha_0,\beta_0,C_{\mathrm{BLM}},\varepsilon_0\).

Fix \(r>0\).  Proposition~\ref{prop:step-004-q-subpower}, at exponent
\(r/(4b)\), and Lemma~\ref{lem:step-004-schedule-subpower}, at
\(r/4\), give
\[
q_\kappa^b\leq L_\kappa^{r/4},\qquad
T_a(q_\kappa,L_\kappa)\leq L_\kappa^{r/4}
\tag{A.84}
\]
on a common tail.  When \(L_\kappa\geq1\),
\[
\begin{aligned}
q_\kappa^b(1+T_a(q_\kappa,L_\kappa))
&\leq
L_\kappa^{r/4}(1+L_\kappa^{r/4})\\
&\leq2L_\kappa^{r/2}.
\end{aligned}
\tag{A.85}
\]
Using the exact ceiling and \(\lceil x\rceil\leq x+1\),
\[
\begin{aligned}
N_\kappa
&=\left\lceil
Hq_\kappa^b(1+T_a(q_\kappa,L_\kappa))
\right\rceil\\
&\leq2HL_\kappa^{r/2}+1\\
&\leq(2H+1)L_\kappa^{r/2}.
\end{aligned}
\tag{A.86}
\]
Eventually \(L_\kappa^{r/2}\geq2H+1\), so
\[
N_\kappa\leq L_\kappa^r.
\tag{A.87}
\]
Thus the fixed power, sum, product, prefactor, additive unit, and
ceiling have each been explicitly controlled.

Because \(N_\kappa\geq1\) and eventually \(L_\kappa>1\),
\[
0\leq\frac{\log N_\kappa}{\log L_\kappa}\leq r
\tag{A.88}
\]
eventually for every \(r>0\), proving the subpower statement.
Taking \(r=1/2\) in (A.81) gives
\[
0\leq\frac{N_\kappa}{L_\kappa}
\leq L_\kappa^{-1/2}
\longrightarrow0,
\tag{A.89}
\]
which proves \(N_\kappa=o(L_\kappa)\) directly.

If \(q_\kappa\leq\bar q\) on a tail, (A.80) and the exact sample
formula yield
\[
N_\kappa
\leq1+H\bar q^b
(1+\bar q^b+u(L_\kappa)).
\tag{A.90}
\]
Using (A.75) at exponent \(1/2\),
\[
\frac{N_\kappa}{L_\kappa}
\leq
\frac{1+H\bar q^b(1+\bar q^b)}{L_\kappa}
+H\bar q^bL_\kappa^{-1/2}
\longrightarrow0.
\tag{A.91}
\]
This retains the bounded-complexity boundary.  Combining
Proposition~\ref{prop:step-004-q-subpower},
Lemma~\ref{lem:step-004-schedule-subpower}, and the present result
gives
\[
T_a(q_\kappa,L_\kappa)=L_\kappa^{o(1)},\qquad
N_\kappa=L_\kappa^{o(1)}=o(L_\kappa).
\tag{A.92}
\]
\end{proof}

\subsection{Uniform Privacy Negligibility}\label{app:privacy-negligibility}

\begin{lemma}[Uniform control of the integer ceiling]
\label{lem:step-005-ceiling}
Under Assumption~\ref{assump:polynomial-global-stability} and
Proposition~\ref{prop:step-004-sample-subpower}, define
\[
x_\kappa=q_\kappa^b,\qquad
u_\kappa=u(L_\kappa),\qquad
A_\kappa=
K_a\varepsilon_0^{-2}
x_\kappa(1+x_\kappa+u_\kappa).
\]
Then, for every \(\kappa\),
\[
A_\kappa\geq6,\qquad
N_\kappa=\lceil A_\kappa\rceil\geq1,\qquad
N_\kappa\leq2A_\kappa.
\tag{A.93}
\]
\end{lemma}

\begin{proof}
The basic setup gives \(q_\kappa\geq1\), hence
\(x_\kappa=q_\kappa^b\geq1\).
Lemma~\ref{lem:step-003-log-correction} gives \(u_\kappa\geq1\), so
\(1+x_\kappa+u_\kappa\geq3\).
Lemma~\ref{lem:step-002-sc-domination} gives \(K_a\geq2\), and
\(\varepsilon_0^{-2}\geq1\).  Therefore
\[
A_\kappa
\geq2\cdot1\cdot1\cdot3=6.
\]
The exact sample formula gives \(N_\kappa=\lceil A_\kappa\rceil\).
Consequently
\[
N_\kappa
\leq A_\kappa+1
\leq A_\kappa+\frac16A_\kappa
\leq2A_\kappa.
\]
The argument is pointwise and imposes no convergence condition on
\(q_\kappa\).
\end{proof}

\begin{lemma}[Two-factor privacy decomposition]
\label{lem:step-005-factorization}
Under Proposition~\ref{prop:step-003-exact-privacy-schedule},
Lemma~\ref{lem:step-005-ceiling}, and a fixed \(\rho>0\), define
\[
C_{\rho,\varepsilon_0,K_a}
=(2K_a\varepsilon_0^{-2})^\rho,
\]
\[
F_\rho(x)=e^{-x}x^\rho(1+x)^\rho,\qquad
H_\rho(u)=e^{-u}(1+u)^\rho.
\]
Then, for every \(\kappa\),
\[
\delta_\kappa N_\kappa^\rho
\leq
C_{\rho,\varepsilon_0,K_a}
F_\rho(x_\kappa)H_\rho(u_\kappa).
\tag{A.94}
\]
\end{lemma}

\begin{proof}
Proposition~\ref{prop:step-003-exact-privacy-schedule} gives
\(\delta_\kappa=e^{-x_\kappa-u_\kappa}\), while
Lemma~\ref{lem:step-005-ceiling} gives
\[
N_\kappa^\rho
\leq
\left[
2K_a\varepsilon_0^{-2}
x_\kappa(1+x_\kappa+u_\kappa)
\right]^\rho.
\]
For nonnegative \(x,u\),
\[
1+x+u\leq(1+x)(1+u),
\tag{A.95}
\]
because the right side is \(1+x+u+xu\).  Since \(\rho>0\), raising
positive quantities to \(\rho\) preserves order, and hence
\[
\begin{aligned}
\delta_\kappa N_\kappa^\rho
&\leq
(2K_a\varepsilon_0^{-2})^\rho
e^{-x_\kappa-u_\kappa}
[x_\kappa(1+x_\kappa+u_\kappa)]^\rho\\
&\leq
(2K_a\varepsilon_0^{-2})^\rho
[e^{-x_\kappa}x_\kappa^\rho(1+x_\kappa)^\rho]
[e^{-u_\kappa}(1+u_\kappa)^\rho].
\end{aligned}
\tag{A.96}
\]
This is (A.94) for the exact \(N_\kappa\) and
\(\delta_\kappa\), with no asymptotic proxy.
\end{proof}

\begin{lemma}[Uniform exponential-polynomial bound]
\label{lem:step-005-x-supremum}
For every fixed \(\rho>0\),
\[
\sup_{x\geq1}F_\rho(x)
\leq
B_\rho
:=
2^\rho e^{-x_\rho}x_\rho^{2\rho}
<\infty,
\qquad
x_\rho=\max\{1,2\rho\}.
\tag{A.97}
\]
\end{lemma}

\begin{proof}
For \(x\geq1\), \(1+x\leq2x\), so
\[
F_\rho(x)\leq2^\rho e^{-x}x^{2\rho}.
\tag{A.98}
\]
Let \(g_\rho(x)=e^{-x}x^{2\rho}\).  On \((0,\infty)\),
\[
g_\rho'(x)
=e^{-x}x^{2\rho-1}(2\rho-x).
\tag{A.99}
\]
If \(2\rho\leq1\), the function is nonincreasing on
\([1,\infty)\), so its maximum there is \(g_\rho(1)\).  If
\(2\rho>1\), it increases on \([1,2\rho]\) and decreases afterward,
so its maximum is \(g_\rho(2\rho)\).  In both cases the maximizer is
\(x_\rho\), and (A.98) gives (A.97).  The constant depends only on
the fixed \(\rho\), not on the sequence.
\end{proof}

\begin{proposition}[Privacy negligibility]
\label{prop:step-005-privacy-negligibility}
Under Assumption~\ref{assump:polynomial-global-stability}, let
\(|C_\kappa|\to\infty\), and let \(N_\kappa,\delta_\kappa\) be the
quantities defined in Section~\ref{sec:prelim}.  For every fixed
\(\rho>0\),
\[
\begin{gathered}
\forall\tau>0\ \exists\kappa_0(\rho,\tau)\
\forall\kappa\geq\kappa_0(\rho,\tau):\\
\delta_\kappa N_\kappa^\rho\leq\tau.
\end{gathered}
\tag{A.100}
\]
Equivalently,
\[
\delta_\kappa N_\kappa^\rho\longrightarrow0
\quad\text{for every fixed }\rho>0.
\tag{A.101}
\]
The conclusion is uniform over all possible behavior of
\(q_\kappa\geq1\), including bounded, constant, oscillatory, and
unbounded sequences.
\end{proposition}

\begin{proof}
Size divergence gives
\[
L_\kappa=\log|C_\kappa|\longrightarrow\infty,
\tag{A.102}
\]
and Lemma~\ref{lem:step-003-log-correction} then gives
\[
u_\kappa=u(L_\kappa)\longrightarrow\infty.
\tag{A.103}
\]
For \(u\geq1\),
\[
\log u
=\int_1^u\frac{dt}{t}
\leq\int_1^u\frac{dt}{\sqrt t}
=2(\sqrt u-1)
\leq2\sqrt u.
\]
Since \(1+u\leq2u\),
\[
\log(1+u)\leq\log2+\log u\leq3\sqrt u.
\tag{A.104}
\]
Whenever \(u\geq36\rho^2\),
\[
\rho\log(1+u)
\leq3\rho\sqrt u
\leq\frac{u}{2},
\qquad
H_\rho(u)\leq e^{-u/2}.
\tag{A.105}
\]
Let
\[
D_\rho=C_{\rho,\varepsilon_0,K_a}B_\rho.
\]
By (A.94) and (A.97),
\[
\delta_\kappa N_\kappa^\rho
\leq D_\rho H_\rho(u_\kappa)
\leq D_\rho e^{-u_\kappa/2}
\tag{A.106}
\]
whenever \(u_\kappa\geq36\rho^2\).
For \(\tau>0\), define
\[
U_{\rho,\tau}
=\max\left\{
1,\ 36\rho^2,\ 2\log\frac{D_\rho}{\tau}
\right\}.
\tag{A.107}
\]
Because \(u_\kappa\to\infty\), eventually
\(u_\kappa\geq U_{\rho,\tau}\), and then
\(D_\rho e^{-u_\kappa/2}\leq\tau\).  If
\(D_\rho/\tau<1\), the logarithmic entry is negative and the desired
inequality is already automatic for \(u_\kappa\geq0\).

The factor \(F_\rho(x_\kappa)\) was bounded uniformly on the entire
domain \(x_\kappa=q_\kappa^b\geq1\).  Thus if \(q_\kappa\) is bounded,
constant, nonconvergent, or oscillatory, the same bound applies while
the independent \(u_\kappa\)-factor vanishes.  If
\(q_\kappa\to\infty\), exponential decay in \(x_\kappa\) can only
improve the estimate.
\end{proof}

\begin{proposition}[Standard approximate-privacy allowance]
\label{prop:step-005-source-allowance}
Under Proposition~\ref{prop:step-005-privacy-negligibility}, for every
fixed \(c>0\) there is \(\kappa_c\) such that
\[
\kappa\geq\kappa_c
\quad\Longrightarrow\quad
\delta_\kappa
\leq
\frac{c}{N_\kappa^2\log(eN_\kappa)}.
\tag{A.108}
\]
\end{proposition}

\begin{proof}
Apply Proposition~\ref{prop:step-005-privacy-negligibility} at the
fixed exponent \(\rho=3\) and tolerance \(c\).  Eventually,
\[
\delta_\kappa N_\kappa^3\leq c.
\tag{A.109}
\]
Lemma~\ref{lem:step-005-ceiling} gives \(N_\kappa\geq1\).  For every
real \(N\geq1\),
\[
\log(eN)=1+\log N\leq N,
\tag{A.110}
\]
because \(f(N)=N-1-\log N\) satisfies \(f(1)=0\) and
\(f'(N)=1-N^{-1}\geq0\).  Hence, on the tail in (A.109),
\[
\delta_\kappa
\leq\frac{c}{N_\kappa^3}
\leq\frac{c}{N_\kappa^2\log(eN_\kappa)}.
\]
The choice \(\rho=3\) is made only after the every-fixed-\(\rho\)
result has been proved, and no growth assumption on \(q_\kappa\) is
used.
\end{proof}

\subsection{Sequence Assembly and the Lower-Bound Contradiction}
\label{app:sequence-assembly}

\begin{proposition}[Sequence-wise exact private learner family]
\label{prop:step-006-sequence-learners}
Under Assumption~\ref{assump:polynomial-global-stability}, fix
\(\alpha_0=\beta_0=1/8\), \(\varepsilon_0\in(0,1)\), and any sequence
\(\{C_\kappa\}\) of nonempty finite binary classes satisfying
\(|C_\kappa|\to\infty\) and (A.57).  There is one constant \(K_a\),
depending only on \(a,\alpha_0,\beta_0,C_{\mathrm{BLM}}\), and for
every \(\kappa\) there is a kernel
\[
M_\kappa:
Z_{X_\kappa}^{N_\kappa}
\rightsquigarrow\mathcal H_{X_\kappa},
\qquad
N_\kappa=N_a(q_\kappa,L_\kappa,\varepsilon_0),
\tag{A.111}
\]
such that:
\begin{enumerate}[label=(\roman*)]
\item \(M_\kappa\) has exact input arity \(N_\kappa\), arbitrary
output, and no representation or computational restriction;
\item for every \(Q\) realizable by \(C_\kappa\),
\[
\Pr_{S\sim Q^{N_\kappa},M_\kappa}
[R_Q(M_\kappa(S))>\alpha_0]\leq\beta_0;
\tag{A.112}
\]
\item \(M_\kappa\) is
\((\varepsilon_0,\delta_\kappa)\)-differentially private, where
\[
\delta_\kappa
=\exp[-q_\kappa^{2a+2}
-\log\log(e^e+L_\kappa)];
\tag{A.113}
\]
\item
\[
\log\frac1{\delta_\kappa}
=q_\kappa^{2a+2}+\log\log(e^e+L_\kappa),
\tag{A.114}
\]
and
\[
\begin{aligned}
N_\kappa
&=
\left\lceil
K_a\varepsilon_0^{-2}q_\kappa^{2a+2}
\left(
1+q_\kappa^{2a+2}
+\log\log(e^e+L_\kappa)
\right)
\right\rceil\\
&\leq
1+K_a\varepsilon_0^{-2}q_\kappa^{2a+2}
\left(
1+q_\kappa^{2a+2}
+\log\log(e^e+L_\kappa)
\right).
\end{aligned}
\tag{A.115}
\]
\end{enumerate}
These conclusions hold at every index, including finite prefixes and
the boundary \(q_\kappa=1\).
\end{proposition}

\begin{proof}
Assumption~\ref{assump:polynomial-global-stability} fixes the same
universal \(a\) before the sequence is considered.
Lemma~\ref{lem:step-002-sc-domination} fixes
\[
K_a=\max\{2,C_{\mathrm{BLM}}D_a\},
\]
independently of
\(\kappa,X_\kappa,C_\kappa,L_\kappa,q_\kappa\), every realizable
distribution, and \(\varepsilon_0\).  For each \(\kappa\), apply
Proposition~\ref{prop:step-002-exact-padding} to
\(C=C_\kappa\), \(X=X_\kappa\), \(q=q_\kappa\),
\(L=L_\kappa\), and the common \(\varepsilon_0\).  It supplies one
kernel \(M_{C_\kappa}\) with exact arity
\(N_a(q_\kappa,L_\kappa,\varepsilon_0)\).  Choose one such kernel at
each index and set
\[
M_\kappa=M_{C_\kappa}.
\tag{A.116}
\]
This indexwise selection does not couple the internal randomness and
does not make the learner depend on \(Q\).

Equation (A.112) and both ordered privacy inequalities are exactly the
conclusions of Proposition~\ref{prop:step-002-exact-padding}.
Proposition~\ref{prop:step-003-exact-privacy-schedule} identifies the
same privacy parameter with (A.113) and proves (A.114).
Substituting \(b=2a+2\) and the definition of \(T_a\) into \(N_a\)
gives the equality in (A.115); \(\lceil x\rceil\leq x+1\) gives its
upper bound.  Since every class is nonempty and finite,
\(L_\kappa\geq0\) and \(q_\kappa\geq1\) at every index, so no
asymptotic tail is needed for this construction.
\end{proof}

\begin{proposition}[Rate certificates for the same learner family]
\label{prop:step-006-rate-certificates}
Under Proposition~\ref{prop:step-006-sequence-learners},
Proposition~\ref{prop:step-004-q-subpower},
Lemma~\ref{lem:step-004-schedule-subpower},
Proposition~\ref{prop:step-004-sample-subpower}, and
Proposition~\ref{prop:step-005-privacy-negligibility}, the exact
quantities in (A.111)--(A.115) satisfy
\[
q_\kappa=L_\kappa^{o(1)},\qquad
T_a(q_\kappa,L_\kappa)=L_\kappa^{o(1)},
\tag{A.117}
\]
\[
N_\kappa=L_\kappa^{o(1)},\qquad
\frac{N_\kappa}{L_\kappa}\longrightarrow0,
\tag{A.118}
\]
and, for every fixed \(\rho>0\),
\[
\delta_\kappa N_\kappa^\rho\longrightarrow0.
\tag{A.119}
\]
The statements include bounded, constant, or oscillatory
\(q_\kappa\), and finite prefixes do not affect the limits.
\end{proposition}

\begin{proof}
Proposition~\ref{prop:step-006-sequence-learners} uses exactly the
quantities \(q_\kappa,L_\kappa,T_a,N_a,\delta_a\) appearing in the
earlier results, so there is no proxy sample size or privacy
parameter.  Proposition~\ref{prop:step-004-q-subpower} and
Lemma~\ref{lem:step-004-schedule-subpower} give (A.117).
Proposition~\ref{prop:step-004-sample-subpower} gives, for every
\(r>0\),
\[
N_\kappa\leq L_\kappa^r
\tag{A.120}
\]
on an eventual tail.  Taking \(r=1/2\) and using
\(L_\kappa\to\infty\) gives
\[
0\leq\frac{N_\kappa}{L_\kappa}
\leq L_\kappa^{-1/2}
\longrightarrow0.
\tag{A.121}
\]
Proposition~\ref{prop:step-005-privacy-negligibility} applies to the
exact pair (A.113), (A.115) and gives (A.119).

If \(q_\kappa\leq\bar q\) on a tail, the direct bound
\[
N_\kappa
\leq
1+K_a\varepsilon_0^{-2}\bar q^b
(1+\bar q^b+u(L_\kappa))
\]
has ratio tending to zero by (A.75), while the proof of
Proposition~\ref{prop:step-005-privacy-negligibility} controls the
\(q_\kappa\)-factor uniformly and lets \(u(L_\kappa)\to\infty\).
This includes \(q_\kappa\equiv1\).
\end{proof}

\begin{lemma}[Privacy relaxation and source-scale eligibility]
\label{lem:step-006-privacy-relaxation}
Under Proposition~\ref{prop:step-006-sequence-learners}, fix
\(\kappa\).  If
\[
\bar\delta(N_\kappa)\geq\delta_\kappa,
\tag{A.122}
\]
then the same \(M_\kappa\) is
\((\varepsilon_0,\bar\delta(N_\kappa))\)-differentially private in
both ordered replacement directions.  Equality is allowed, and no
reverse implication is claimed when
\(\bar\delta(N_\kappa)<\delta_\kappa\).

Moreover, for every fixed \(c>0\), the same learner is eventually
\[
\left(
\varepsilon_0,
\frac{c}{N_\kappa^2\log(eN_\kappa)}
\right)\text{-differentially private}.
\tag{A.123}
\]
\end{lemma}

\begin{proof}
Let \(S,S'\in Z_{X_\kappa}^{N_\kappa}\) be replacement-adjacent and
let \(E\subseteq\mathcal H_{X_\kappa}\) be measurable.
Proposition~\ref{prop:step-006-sequence-learners} gives
\[
\Pr[M_\kappa(S)\in E]
\leq e^{\varepsilon_0}\Pr[M_\kappa(S')\in E]
+\delta_\kappa.
\]
Using (A.122) in the valid direction gives
\[
\Pr[M_\kappa(S)\in E]
\leq e^{\varepsilon_0}\Pr[M_\kappa(S')\in E]
+\bar\delta(N_\kappa).
\tag{A.124}
\]
The reverse ordered inequality similarly becomes
\[
\Pr[M_\kappa(S')\in E]
\leq e^{\varepsilon_0}\Pr[M_\kappa(S)\in E]
+\bar\delta(N_\kappa).
\tag{A.125}
\]
Thus only the allowed additive slack has increased; the kernel,
sample count, output space, risk, and adjacency relation are
unchanged.

Proposition~\ref{prop:step-005-source-allowance} states that, for every
fixed \(c>0\), eventually
\[
\delta_\kappa
\leq\frac{c}{N_\kappa^2\log(eN_\kappa)}.
\]
Substitution into (A.124)--(A.125) gives (A.123).  At the exact
schedule no tail is needed; if a comparison holds only eventually,
the relaxed conclusion holds on exactly that tail.
\end{proof}

\begin{proposition}[No eventual linear lower bound]
\label{prop:step-006-omega-contradiction}
Under Proposition~\ref{prop:step-006-sequence-learners},
Proposition~\ref{prop:step-006-rate-certificates}, and
Lemma~\ref{lem:step-006-privacy-relaxation}, consider either the exact
schedule \(\delta_\kappa\), or a schedule \(\bar\delta\) for which
\[
\bar\delta(N_\kappa)\geq\delta_\kappa
\quad\text{eventually}.
\tag{A.126}
\]
There do not exist \(c_*>0\) and \(\kappa_*\) such that, for every
\(\kappa\geq\kappa_*\), every arbitrary-output, computationally
unrestricted, distribution-free realizable
\((\alpha_0,\beta_0)\)-PAC learner for \(C_\kappa\) at the relevant
privacy schedule must use at least
\[
c_*L_\kappa
\tag{A.127}
\]
records.
\end{proposition}

\begin{proof}
Fix arbitrary \(c_*>0\) and \(\kappa_*\).  In the exact-schedule case
there is no privacy tail; otherwise let
\(\kappa_{\mathrm{priv}}\) be a tail index for (A.126).
Proposition~\ref{prop:step-006-rate-certificates} gives
\(N_\kappa/L_\kappa\to0\), so there is
\(\kappa_{\mathrm{rate}}(c_*)\) such that
\[
N_\kappa<c_*L_\kappa
\qquad
(\kappa\geq\kappa_{\mathrm{rate}}(c_*)).
\tag{A.128}
\]
The ratio is well-defined there because \(L_\kappa\to\infty\).
Choose
\[
\kappa\geq
\max\{\kappa_*,\kappa_{\mathrm{priv}},
\kappa_{\mathrm{rate}}(c_*)\}.
\tag{A.129}
\]
Proposition~\ref{prop:step-006-sequence-learners} supplies an eligible
learner \(M_\kappa\) using exactly \(N_\kappa\) records.  In the
comparison-schedule case,
Lemma~\ref{lem:step-006-privacy-relaxation} makes the same learner
eligible.  The hypothetical lower bound would give
\(N_\kappa\geq c_*L_\kappa\), contradicting (A.128).
Since \(c_*\) and the proposed tail were arbitrary, no positive
eventual linear lower-bound constant exists.  The standard allowance
follows by taking
\(\bar\delta(N_\kappa)=
c/[N_\kappa^2\log(eN_\kappa)]\) and using (A.123).
\end{proof}

\subsection{Proof of the Main Theorem}\label{app:main-proof}

\begin{proof}[Proof of Theorem~\ref{thm:main}]
Assumption~\ref{assump:polynomial-global-stability} and
Proposition~\ref{prop:step-001-binary-private-learner} give the
arbitrary-output stable-selection learner and its exact two-term raw
sample bound.  Lemma~\ref{lem:step-002-sc-domination} proves the
displayed sample specialization and fixes a common \(K_a\) with only
the allowed dependence; Proposition~\ref{prop:step-002-exact-padding}
turns the learner into an exact-\(N_a\) kernel without changing
utility or privacy.  Proposition~\ref{prop:step-003-exact-privacy-schedule}
proves the displayed value of \(\delta_\kappa\) and its exact
logarithm.

For an arbitrary admissible sequence,
Proposition~\ref{prop:step-004-q-subpower},
Lemma~\ref{lem:step-004-schedule-subpower}, and
Proposition~\ref{prop:step-004-sample-subpower} prove
\[
q_\kappa=L_\kappa^{o(1)},\qquad
T_a(q_\kappa,L_\kappa)=L_\kappa^{o(1)},\qquad
N_\kappa=L_\kappa^{o(1)}=o(L_\kappa).
\]
Proposition~\ref{prop:step-005-privacy-negligibility} proves
\(\delta_\kappa N_\kappa^\rho\to0\) for every fixed \(\rho>0\), and
Proposition~\ref{prop:step-005-source-allowance} gives the stated
standard approximate-privacy comparison.

Proposition~\ref{prop:step-006-sequence-learners} attaches all
pointwise guarantees to one learner \(M_\kappa\) for each index, and
Proposition~\ref{prop:step-006-rate-certificates} attaches the exact
sequence rates to those same learners.  Lemma~\ref{lem:step-006-privacy-relaxation}
proves privacy monotonicity only toward a larger allowed additive
parameter, in both ordered adjacency directions.
Proposition~\ref{prop:step-006-omega-contradiction} then combines the
valid learner with \(N_\kappa/L_\kappa\to0\) and rules out every
eventual positive linear lower bound at the exact, standard, or any
eventually weaker privacy allowance.  The argument never removes
Assumption~\ref{assump:polynomial-global-stability}; hence the result
is conditional exactly as stated.
\end{proof}
